% Source: book pp. 308--318 (PDF 322--332). No solutions.
% numbered displays (Axler's 8.32, ...): see 8A.tex
\DeclareRobustCommand\axeqnum[1]{\refstepcounter{theorem}\label{ax:#1}}
\section{Generalized Eigenspace Decomposition}

\subsection{Generalized Eigenspaces}

\begin{definition}[Generalized eigenspace, $G(\lambda,T)$]\label{ax:8.19}
Suppose $T\in\Lin(V)$ and $\lambda\in\F$. The \emph{generalized eigenspace} of
$T$ corresponding to $\lambda$, denoted by $G(\lambda,T)$, is defined by
\[ G(\lambda,T) = \{v\in V : (T-\lambda I)^kv=0 \text{ for some positive
   integer } k\}. \]
Thus $G(\lambda,T)$ is the set of generalized eigenvectors of $T$ corresponding
to $\lambda$, along with the $0$ vector.
\end{definition}

Because every eigenvector of $T$ is a generalized eigenvector of $T$ (take
$k=1$ in the definition of generalized eigenvector), each eigenspace is
contained in the corresponding generalized eigenspace. In other words, if
$T\in\Lin(V)$ and $\lambda\in\F$, then $E(\lambda,T)\subseteq G(\lambda,T)$.

The next result implies that if $T\in\Lin(V)$ and $\lambda\in\F$, then the
generalized eigen\-space $G(\lambda,T)$ is a subspace of $V$ (because the null
space of each linear map on $V$ is a subspace of $V$).

\begin{result}[Description of generalized eigenspaces]\label{ax:8.20}
Suppose $T\in\Lin(V)$ and $\lambda\in\F$. Then
$G(\lambda,T)=\nul(T-\lambda I)^{\dim V}$.
\end{result}

\begin{proof}
Suppose $v\in\nul(T-\lambda I)^{\dim V}$. The definitions imply
$v\in G(\lambda,T)$. Thus $G(\lambda,T)\supseteq\nul(T-\lambda I)^{\dim V}$.

Conversely, suppose $v\in G(\lambda,T)$. Thus there is a positive integer $k$
such that $v\in\nul(T-\lambda I)^k$. From \ref{ax:8.1} and \ref{ax:8.3} (with
$T-\lambda I$ replacing $T$), we get $v\in\nul(T-\lambda I)^{\dim V}$. Thus
$G(\lambda,T)\subseteq\nul(T-\lambda I)^{\dim V}$, completing the proof.
\end{proof}

\begin{example}[Generalized eigenspaces of an operator on $\C^3$]\label{ax:8.21}
Define $T\in\Lin(\C^3)$ by
\[ T(z_1,z_2,z_3) = (4z_2,0,5z_3). \]

In Example \ref{ax:8.10}, we saw that the eigenvalues of $T$ are $0$ and $5$,
and we found the corresponding sets of generalized eigenvectors. Taking the
union of those sets with $\{0\}$, we have
\[ G(0,T) = \{(z_1,z_2,0) : z_1,z_2\in\C\} \quad\text{and}\quad
   G(5,T) = \{(0,0,z_3) : z_3\in\C\}. \]
Note that $\C^3=G(0,T)\oplus G(5,T)$.
\end{example}

In Example \ref{ax:8.21}, the domain space $\C^3$ is the direct sum of the
generalized eigenspaces of the operator $T$ in that example. Our next result
shows that this behavior holds in general. Specifically, the following major
result shows that if $\F=\C$ and $T\in\Lin(V)$, then $V$ is the direct sum of
the generalized eigenspaces of $T$, each of which is invariant under $T$ and on
which $T$ is a nilpotent operator plus a scalar multiple of the identity. Thus
the next result achieves our goal of decomposing $V$ into invariant subspaces on
which $T$ has a known behavior.

As we will see, the proof follows from putting together what we have learned
about generalized eigenspaces and then using our result that for each operator
$T\in\Lin(V)$, there exists a basis of $V$ consisting of generalized
eigenvectors of $T$.

\begin{result}[Generalized eigenspace decomposition]\label{ax:8.22}\emergencystretch=1.5em
Suppose $\F=\C$ and $T\in\Lin(V)$. Let $\lambda_1,\dots,\lambda_m$ be the
distinct eigenvalues of $T$. Then
\begin{parts}
\item $G(\lambda_k,T)$ is invariant under $T$ for each $k=1,\dots,m$;
\item $(T-\lambda_kI)|_{G(\lambda_k,T)}$ is nilpotent for each
  $k=1,\dots,m$;
\item $V=G(\lambda_1,T)\oplus\cdots\oplus G(\lambda_m,T)$.
\end{parts}
\end{result}

\begin{proof}\leavevmode
\begin{parts}
\item Suppose $k\in\{1,\dots,m\}$. Then \ref{ax:8.20} shows that
  \[ G(\lambda_k,T) = \nul(T-\lambda_kI)^{\dim V}. \]
  Thus \ref{ax:5.18}, with $p(z)=(z-\lambda_k)^{\dim V}$, implies that
  $G(\lambda_k,T)$ is invariant under $T$, proving (a).
\item Suppose $k\in\{1,\dots,m\}$. If $v\in G(\lambda_k,T)$, then
  $(T-\lambda_kI)^{\dim V}v=0$ (by \ref{ax:8.20}). Thus
  $\bigl((T-\lambda_kI)|_{G(\lambda_k,T)}\bigr)^{\dim V}=0$. Hence
  $(T-\lambda_kI)|_{G(\lambda_k,T)}$ is nilpotent, proving (b).
\item To show that $G(\lambda_1,T)+\cdots+G(\lambda_m,T)$ is a direct sum,
  suppose
  \[ v_1+\cdots+v_m = 0, \]
  where each $v_k$ is in $G(\lambda_k,T)$. Because generalized eigenvectors of
  $T$ corresponding to distinct eigenvalues are linearly independent (by
  \ref{ax:8.12}), this implies that each $v_k$ equals $0$. Thus
  $G(\lambda_1,T)+\cdots+G(\lambda_m,T)$ is a direct sum (by \ref{ax:1.45}).

  Finally, each vector in $V$ can be written as a finite sum of generalized
  eigenvectors of $T$ (by \ref{ax:8.9}). Thus
  \[ V = G(\lambda_1,T)\oplus\cdots\oplus G(\lambda_m,T), \]
  proving (c).\qedhere
\end{parts}
\end{proof}

For the analogous result when $\F=\R$, see Exercise 8.

\subsection{Multiplicity of an Eigenvalue}

If $V$ is a complex vector space and $T\in\Lin(V)$, then the decomposition of
$V$ provided by the generalized eigenspace decomposition (\ref{ax:8.22}) can be
a powerful tool. The dimensions of the subspaces involved in this decomposition
are sufficiently important to get a name, which is given in the next
definition.

\begin{definition}[Multiplicity]\label{ax:8.23}\leavevmode
\begin{itemize}
\item Suppose $T\in\Lin(V)$. The \emph{multiplicity} of an eigenvalue $\lambda$
  of $T$ is defined to be the dimension of the corresponding generalized
  eigenspace $G(\lambda,T)$.
\item In other words, the multiplicity of an eigenvalue $\lambda$ of $T$ equals
  \[ \dim\nul(T-\lambda I)^{\dim V}. \]
\end{itemize}
\end{definition}

The second bullet point above holds because
$G(\lambda,T)=\nul(T-\lambda I)^{\dim V}$ (see \ref{ax:8.20}).

\begin{example}[Multiplicity of each eigenvalue of an operator]\label{ax:8.24}
Suppose $T\in\Lin(\C^3)$ is defined by
\[ T(z_1,z_2,z_3) = (6z_1+3z_2+4z_3,6z_2+2z_3,7z_3). \]
The matrix of $T$ (with respect to the standard basis) is
\[ \begin{pmatrix} 6 & 3 & 4\\ 0 & 6 & 2\\ 0 & 0 & 7 \end{pmatrix}. \]
The eigenvalues of $T$ are the diagonal entries $6$ and $7$, as follows from
\ref{ax:5.41}. You can verify that the generalized eigenspaces of $T$ are as
follows:
\[ G(6,T) = \spn\bigl((1,0,0),(0,1,0)\bigr) \quad\text{and}\quad
   G(7,T) = \spn\bigl((10,2,1)\bigr). \]

Thus the eigenvalue $6$ has multiplicity $2$ and the eigenvalue $7$ has
multiplicity $1$. The direct sum $\C^3=G(6,T)\oplus G(7,T)$ is the generalized
eigenspace decomposition promised by \ref{ax:8.22}. A basis of $\C^3$
consisting of generalized eigenvectors of $T$, as promised by \ref{ax:8.9}, is
$(1,0,0),(0,1,0),(10,2,1)$. There does not exist a basis of $\C^3$ consisting
of eigenvectors of this operator.
\end{example}

\marginnote[-24mm]{In this example, the multiplicity of each eigenvalue equals the
number of times that eigenvalue appears on the diagonal of an upper-triangular
matrix representing the operator. This behavior always happens, as we will see
in \ref{ax:8.31}.}%
In the example above, the sum of the multiplicities of the eigenvalues of $T$
equals $3$, which is the dimension of the domain of $T$. The next result shows
that this holds for all operators on finite-dimensional complex vector spaces.

\begin{result}[Sum of the multiplicities equals $\dim V$]\label{ax:8.25}
Suppose $\F=\C$ and $T\in\Lin(V)$. Then the sum of the multiplicities of all
eigenvalues of $T$ equals $\dim V$.
\end{result}

\begin{proof}
The desired result follows from the generalized eigenspace decomposition
(\ref{ax:8.22}) and the formula for the dimension of a direct sum (see
\ref{ax:3.94}).
\end{proof}

The terms \emph{algebraic multiplicity} and \emph{geometric multiplicity} are
used in some books. In case you encounter this terminology, be aware that the
algebraic multiplicity is the same as the multiplicity defined here and the
geometric multiplicity is the dimension of the corresponding eigenspace. In
other words, if $T\in\Lin(V)$ and $\lambda$ is an eigenvalue of $T$, then
\begin{align*}
\text{algebraic multiplicity of } \lambda
  &= \dim\nul(T-\lambda I)^{\dim V} = \dim G(\lambda,T),\\
\text{geometric multiplicity of } \lambda
  &= \dim\nul(T-\lambda I) = \dim E(\lambda,T).
\end{align*}
Note that as defined above, the algebraic multiplicity also has a geometric
meaning as the dimension of a certain null space. The definition of
multiplicity given here is cleaner than the traditional definition that
involves determinants; \ref{ax:9.62} implies that these definitions are
equivalent.

If $V$ is an inner product space, $T\in\Lin(V)$ is normal, and $\lambda$ is an
eigenvalue of $T$, then the algebraic multiplicity of $\lambda$ equals the
geometric multiplicity of $\lambda$, as can be seen from applying Exercise 27
in Section 7A to the normal operator $T-\lambda I$. As a special case, the
singular values of $S\in\Lin(V,W)$ (here $V$ and $W$ are both
finite-dimensional inner product spaces) depend on the multiplicities (either
algebraic or geometric) of the eigenvalues of the self-adjoint operator
$S^*S$.

The next definition associates a monic polynomial with each operator on a
finite-dimensional complex vector space.

\begin{definition}[Characteristic polynomial]\label{ax:8.26}
Suppose $\F=\C$ and $T\in\Lin(V)$. Let $\lambda_1,\dots,\lambda_m$ denote the
distinct eigenvalues of $T$, with multiplicities $d_1,\dots,d_m$. The
polynomial
\[ (z-\lambda_1)^{d_1}\cdots(z-\lambda_m)^{d_m} \]
is called the \emph{characteristic polynomial} of $T$.
\end{definition}

\begin{example}[The characteristic polynomial of an operator]\label{ax:8.27}
Suppose $T\in\Lin(\C^3)$ is defined as in Example \ref{ax:8.24}. Because the
eigenvalues of $T$ are $6$, with multiplicity $2$, and $7$, with multiplicity
$1$, we see that the characteristic polynomial of $T$ is $(z-6)^2(z-7)$.
\end{example}

\begin{result}[Degree and zeros of characteristic polynomial]\label{ax:8.28}
Suppose $\F=\C$ and $T\in\Lin(V)$. Then
\begin{parts}
\item the characteristic polynomial of $T$ has degree $\dim V$;
\item the zeros of the characteristic polynomial of $T$ are the eigenvalues of
  $T$.
\end{parts}
\end{result}

\begin{proof}
Our result about the sum of the multiplicities (\ref{ax:8.25}) implies (a). The
definition of the characteristic polynomial implies (b).
\end{proof}

Most texts define the characteristic polynomial using determinants (the two
definitions are equivalent by \ref{ax:9.62}). The approach taken here, which is
considerably simpler, leads to the following nice proof of the Cayley--Hamilton
theorem.

\begin{result}[Cayley--Hamilton theorem]\label{ax:8.29}
Suppose $\F=\C$, $T\in\Lin(V)$, and $q$ is the characteristic polynomial of
$T$. Then $q(T)=0$.
\end{result}

\begin{proof}
Let $\lambda_1,\dots,\lambda_m$ be the distinct eigenvalues of $T$, and let
$d_k=\dim G(\lambda_k,T)$. For each $k\in\{1,\dots,m\}$, we know that
$(T-\lambda_kI)|_{G(\lambda_k,T)}$ is nilpotent. Thus we have
\[ (T-\lambda_kI)^{d_k}|_{G(\lambda_k,T)} = 0 \]
(by \ref{ax:8.16}) for each $k\in\{1,\dots,m\}$.

\marginnote{Arthur Cayley (1821--1895) published three mathematics
papers before completing his undergraduate degree.}%
The generalized eigenspace decomposition (\ref{ax:8.22}) states that every
vector in $V$ is a sum of vectors in $G(\lambda_1,T),\dots,G(\lambda_m,T)$.
Thus to prove that $q(T)=0$, we only need to show that
$q(T)|_{G(\lambda_k,T)}=0$ for each $k$.

Fix $k\in\{1,\dots,m\}$. We have
\[ q(T) = (T-\lambda_1I)^{d_1}\cdots(T-\lambda_mI)^{d_m}. \]
The operators on the right side of the equation above all commute, so we can
move the factor $(T-\lambda_kI)^{d_k}$ to be the last term in the expression on
the right. Because $(T-\lambda_kI)^{d_k}|_{G(\lambda_k,T)}=0$, we have
$q(T)|_{G(\lambda_k,T)}=0$, as desired.
\end{proof}

The next result implies that if the minimal polynomial of an operator
$T\in\Lin(V)$ has degree $\dim V$ (as happens almost always---see the
paragraphs following \ref{ax:5.24}), then the characteristic polynomial of $T$
equals the minimal polynomial of $T$.

\begin{result}[Characteristic polynomial is a multiple of minimal
polynomial]\label{ax:8.30}
Suppose $\F=\C$ and $T\in\Lin(V)$. Then the characteristic polynomial of $T$ is
a polynomial multiple of the minimal polynomial of $T$.
\end{result}

\begin{proof}
The desired result follows immediately from the Cayley--Hamilton theorem
(\ref{ax:8.29}) and \ref{ax:5.29}.
\end{proof}

Now we can prove that the result suggested by Example \ref{ax:8.24} holds for
all operators on finite-dimensional complex vector spaces.

\begin{result}[Multiplicity of an eigenvalue equals number of times on
diagonal]\label{ax:8.31}
Suppose $\F=\C$ and $T\in\Lin(V)$. Suppose $v_1,\dots,v_n$ is a basis of $V$
such that $\Mat\bigl(T,(v_1,\dots,v_n)\bigr)$ is upper triangular. Then the
number of times that each eigenvalue $\lambda$ of $T$ appears on the diagonal of
$\Mat\bigl(T,(v_1,\dots,v_n)\bigr)$ equals the multiplicity of $\lambda$ as an
eigenvalue of $T$.
\end{result}

\begin{proof}
Let $A=\Mat\bigl(T,(v_1,\dots,v_n)\bigr)$. Thus $A$ is an upper-triangular
matrix. Let $\lambda_1,\dots,\lambda_n$ denote the entries on the diagonal of
$A$. Thus for each $k\in\{1,\dots,n\}$, we have
\[ Tv_k = u_k+\lambda_kv_k \]
for some $u_k\in\spn(v_1,\dots,v_{k-1})$. Hence if $k\in\{1,\dots,n\}$ and
$\lambda_k\ne0$, then $Tv_k$ is not a linear combination of
$Tv_1,\dots,Tv_{k-1}$. The linear dependence lemma (\ref{ax:2.19}) now implies
that the list of those $Tv_k$ such that $\lambda_k\ne0$ is linearly
independent.

Let $d$ denote the number of indices $k\in\{1,\dots,n\}$ such that
$\lambda_k=0$. The conclusion of the previous paragraph implies that
\[ \dim\range T \ge n-d. \]
Because $n=\dim V=\dim\nul T+\dim\range T$, the inequality above implies
that\axeqnum{8.32}%
\[ \dim\nul T \le d. \tag*{\thetheorem} \]

The matrix of the operator $T^n$ with respect to the basis $v_1,\dots,v_n$ is
the upper-triangular matrix $A^n$, which has diagonal entries
$\lambda_1^n,\dots,\lambda_n^n$ $\bigl[$see Exercise 2(b) in Section 5C$\bigr]$.
Because $\lambda_k^n=0$ if and only if $\lambda_k=0$, the number of times that
$0$ appears on the diagonal of $A^n$ equals $d$. Thus applying \ref{ax:8.32}
with $T$ replaced with $T^n$, we have\axeqnum{8.33}%
\[ \dim\nul T^n \le d. \tag*{\thetheorem} \]

For $\lambda$ an eigenvalue of $T$, let $m_\lambda$ denote the multiplicity of
$\lambda$ as an eigenvalue of $T$ and let $d_\lambda$ denote the number of
times that $\lambda$ appears on the diagonal of $A$. Replacing $T$ in
\ref{ax:8.33} with $T-\lambda I$, we see that\axeqnum{8.34}%
\[ m_\lambda \le d_\lambda \tag*{\thetheorem} \]
for each eigenvalue $\lambda$ of $T$. The sum of the multiplicities $m_\lambda$
over all eigenvalues $\lambda$ of $T$ equals $n$, the dimension of $V$ (by
\ref{ax:8.25}). The sum of the numbers $d_\lambda$ over all eigenvalues
$\lambda$ of $T$ also equals $n$, because the diagonal of $A$ has length $n$.

Thus summing both sides of \ref{ax:8.34} over all eigenvalues $\lambda$ of $T$
produces an equality. Hence \ref{ax:8.34} must actually be an equality for each
eigenvalue $\lambda$ of $T$. Thus the multiplicity of $\lambda$ as an
eigenvalue of $T$ equals the number of times that $\lambda$ appears on the
diagonal of $A$, as desired.
\end{proof}

\subsection{Block Diagonal Matrices}

\marginnote{Often we can understand a matrix better by thinking of it as
composed of smaller matrices.}%
To interpret our results in matrix form, we make the following definition,
generalizing the notion of a diagonal matrix. If each matrix $A_k$ in the
definition below is a $1$-by-$1$ matrix, then we actually have a diagonal
matrix.

\begin{definition}[Block diagonal matrix]\label{ax:8.35}
A \emph{block diagonal matrix} is a square matrix of the form
\[ \begin{pmatrix} A_1 & & 0\\ & \ddots & \\ 0 & & A_m \end{pmatrix}, \]
where $A_1,\dots,A_m$ are square matrices lying along the diagonal and all
other entries of the matrix equal $0$.
\end{definition}

\begin{example}[A block diagonal matrix]\label{ax:8.36}
The $5$-by-$5$ matrix
\[ A = \begin{pmatrix}
  \begin{pmatrix} 4 \end{pmatrix} & \begin{matrix} 0 & 0 \end{matrix}
    & \begin{matrix} 0 & 0 \end{matrix}\\
  \begin{matrix} 0\\ 0 \end{matrix}
    & \begin{pmatrix} 2 & -3\\ 0 & 2 \end{pmatrix}
    & \begin{matrix} 0 & 0\\ 0 & 0 \end{matrix}\\
  \begin{matrix} 0\\ 0 \end{matrix}
    & \begin{matrix} 0 & 0\\ 0 & 0 \end{matrix}
    & \begin{pmatrix} 1 & 7\\ 0 & 1 \end{pmatrix}
\end{pmatrix} \]
is a block diagonal matrix with
\[ A = \begin{pmatrix} A_1 & & 0\\ & A_2 & \\ 0 & & A_3 \end{pmatrix}, \]
where
\[ A_1 = \begin{pmatrix} 4 \end{pmatrix}, \quad
   A_2 = \begin{pmatrix} 2 & -3\\ 0 & 2 \end{pmatrix}, \quad
   A_3 = \begin{pmatrix} 1 & 7\\ 0 & 1 \end{pmatrix}. \]
Here the inner matrices in the $5$-by-$5$ matrix above are blocked off to show
how we can think of it as a block diagonal matrix.
\end{example}

Note that in the example above, each of $A_1,A_2,A_3$ is an upper-triangular
matrix whose diagonal entries are all equal. The next result shows that with
respect to an appropriate basis, every operator on a finite-dimensional complex
vector space has a matrix of this form. Note that this result gives us many
more zeros in the matrix than are needed to make it upper triangular.

\begin{result}[Block diagonal matrix with upper-triangular blocks]\label{ax:8.37}
Suppose $\F=\C$ and $T\in\Lin(V)$. Let $\lambda_1,\dots,\lambda_m$ be the
distinct eigenvalues of $T$, with multiplicities $d_1,\dots,d_m$. Then there is
a basis of $V$ with respect to which $T$ has a block diagonal matrix of the
form
\[ \begin{pmatrix} A_1 & & 0\\ & \ddots & \\ 0 & & A_m \end{pmatrix}, \]
where each $A_k$ is a $d_k$-by-$d_k$ upper-triangular matrix of the form
\[ A_k = \begin{pmatrix} \lambda_k & & *\\ & \ddots & \\ 0 & & \lambda_k
   \end{pmatrix}. \]
\end{result}

\begin{proof}
Each $(T-\lambda_kI)|_{G(\lambda_k,T)}$ is nilpotent (see \ref{ax:8.22}). For
each $k$, choose a basis of $G(\lambda_k,T)$, which is a vector space of
dimension $d_k$, such that the matrix of $(T-\lambda_kI)|_{G(\lambda_k,T)}$
with respect to this basis is as in \ref{ax:8.18}(c). Thus with respect to this
basis, the matrix of $T|_{G(\lambda_k,T)}$, which equals
$(T-\lambda_kI)|_{G(\lambda_k,T)}+\lambda_kI|_{G(\lambda_k,T)}$, looks like the
desired form shown above for $A_k$.

The generalized eigenspace decomposition (\ref{ax:8.22}) shows that putting
together the bases of the $G(\lambda_k,T)$'s chosen above gives a basis of $V$.
The matrix of $T$ with respect to this basis has the desired form.
\end{proof}

\begin{example}[Block diagonal matrix via generalized eigenvectors]\label{ax:8.38}
Let $T\in\Lin(\C^3)$ be defined by
$T(z_1,z_2,z_3)=(6z_1+3z_2+4z_3,6z_2+2z_3,7z_3)$. The matrix of $T$ (with
respect to the standard basis) is
\[ \begin{pmatrix} 6 & 3 & 4\\ 0 & 6 & 2\\ 0 & 0 & 7 \end{pmatrix}, \]
which is an upper-triangular matrix but is not of the form promised by
\ref{ax:8.37}.

As we saw in Example \ref{ax:8.24}, the eigenvalues of $T$ are $6$ and $7$;
also,
\[ G(6,T) = \spn\bigl((1,0,0),(0,1,0)\bigr) \quad\text{and}\quad
   G(7,T) = \spn\bigl((10,2,1)\bigr). \]
We also saw that a basis of $\C^3$ consisting of generalized eigenvectors of
$T$ is
\[ (1,0,0),(0,1,0),(10,2,1). \]
The matrix of $T$ with respect to this basis is
\[ \begin{pmatrix}
  \begin{pmatrix} 6 & 3\\ 0 & 6 \end{pmatrix} & \begin{matrix} 0\\ 0 \end{matrix}\\
  \begin{matrix} 0 & 0 \end{matrix} & \begin{pmatrix} 7 \end{pmatrix}
\end{pmatrix}, \]
which is a matrix of the block diagonal form promised by \ref{ax:8.37}.
\end{example}

% ------------------------------------------------------------------ exercises
\begin{exc}

\begin{exercise}
Define $T\in\Lin(\C^2)$ by $T(w,z)=(-z,w)$. Find the generalized eigenspaces
corresponding to the distinct eigenvalues of $T$.
\end{exercise}

\begin{exercise}
Suppose $T\in\Lin(V)$ is invertible. Prove that
$G(\lambda,T)=G(1/\lambda,T^{-1})$ for every $\lambda\in\F$
with $\lambda\ne0$.
\end{exercise}

\begin{exercise}
Suppose $T\in\Lin(V)$. Suppose $S\in\Lin(V)$ is invertible. Prove that $T$ and
$S^{-1}TS$ have the same eigenvalues with the same multiplicities.
\end{exercise}

\begin{exercise}
Suppose $\dim V\ge2$ and $T\in\Lin(V)$ is such that
$\nul T^{\dim V-2}\ne\nul T^{\dim V-1}$. Prove that $T$ has at most two
distinct eigenvalues.
\end{exercise}

\begin{exercise}
Suppose $T\in\Lin(V)$ and $3$ and $8$ are eigenvalues of $T$. Let $n=\dim V$.
Prove that $V=(\nul T^{n-2})\oplus(\range T^{n-2})$.
\end{exercise}

\begin{exercise}
Suppose $T\in\Lin(V)$ and $\lambda$ is an eigenvalue of $T$. Explain why the
exponent of $z-\lambda$ in the factorization of the minimal polynomial of $T$
is the smallest positive integer $m$ such that
$(T-\lambda I)^m|_{G(\lambda,T)}=0$.
\end{exercise}

\begin{exercise}
Suppose $T\in\Lin(V)$ and $\lambda$ is an eigenvalue of $T$ with multiplicity
$d$. Prove that $G(\lambda,T)=\nul(T-\lambda I)^d$.
\exnote{If $d<\dim V$, then this exercise improves \ref{ax:8.20}.}
\end{exercise}

\begin{exercise}
Suppose $T\in\Lin(V)$ and $\lambda_1,\dots,\lambda_m$ are the distinct
eigenvalues of $T$. Prove that
\[ V = G(\lambda_1,T)\oplus\cdots\oplus G(\lambda_m,T) \]
if and only if the minimal polynomial of $T$ equals
$(z-\lambda_1)^{k_1}\cdots(z-\lambda_m)^{k_m}$ for some positive integers
$k_1,\dots,k_m$.
\exnote{The case $\F=\C$ follows immediately from
\ref{ax:5.27}\textup{(b)} and the generalized eigenspace decomposition
\textup{(\ref{ax:8.22})}; thus this exercise is interesting only when
$\F=\R$.}
\end{exercise}

\begin{exercise}
Suppose $\F=\C$ and $T\in\Lin(V)$. Prove that there exist $D,N\in\Lin(V)$ such
that $T=D+N$, the operator $D$ is diagonalizable, $N$ is nilpotent, and
$DN=ND$.
\end{exercise}

\begin{exercise}
Suppose $V$ is a complex inner product space, $e_1,\dots,e_n$ is an orthonormal
basis of $V$, and $T\in\Lin(V)$. Let $\lambda_1,\dots,\lambda_n$ be the
eigenvalues of $T$, each included as many times as its multiplicity. Prove that
\[ \abs{\lambda_1}^2+\cdots+\abs{\lambda_n}^2
   \le \norm{Te_1}^2+\cdots+\norm{Te_n}^2. \]
\exnote{See the comment after Exercise 5 in Section 7A.}
\end{exercise}

\begin{exercise}
Give an example of an operator on $\C^4$ whose characteristic polynomial equals
$(z-7)^2(z-8)^2$.
\end{exercise}

\begin{exercise}
Give an example of an operator on $\C^4$ whose characteristic polynomial equals
$(z-1)(z-5)^3$ and whose minimal polynomial equals $(z-1)(z-5)^2$.
\end{exercise}

\begin{exercise}
Give an example of an operator on $\C^4$ whose characteristic and minimal
polynomials both equal $z(z-1)^2(z-3)$.
\end{exercise}

\begin{exercise}
Give an example of an operator on $\C^4$ whose characteristic polynomial equals
$z(z-1)^2(z-3)$ and whose minimal polynomial equals $z(z-1)(z-3)$.
\end{exercise}

\begin{exercise}
Let $T$ be the operator on $\C^4$ defined by
$T(z_1,z_2,z_3,z_4)=(0,z_1,z_2,z_3)$. Find the characteristic polynomial and
the minimal polynomial of $T$.
\end{exercise}

\begin{exercise}
Let $T$ be the operator on $\C^6$ defined by
\[ T(z_1,z_2,z_3,z_4,z_5,z_6) = (0,z_1,z_2,0,z_4,0). \]
Find the characteristic polynomial and the minimal polynomial of $T$.
\end{exercise}

\begin{exercise}
Suppose $\F=\C$ and $P\in\Lin(V)$ is such that $P^2=P$. Prove that the
characteristic polynomial of $P$ is $z^m(z-1)^n$, where $m=\dim\nul P$ and
$n=\dim\range P$.
\end{exercise}

\begin{exercise}
Suppose $T\in\Lin(V)$ and $\lambda$ is an eigenvalue of $T$. Explain why the
following four numbers equal each other.
\begin{parts}
\item The exponent of $z-\lambda$ in the factorization of the minimal
  polynomial of $T$.
\item The smallest positive integer $m$ such that
  $(T-\lambda I)^m|_{G(\lambda,T)}=0$.
\item The smallest positive integer $m$ such that
  \[ \nul(T-\lambda I)^m = \nul(T-\lambda I)^{m+1}. \]
\item The smallest positive integer $m$ such that
  \[ \range(T-\lambda I)^m = \range(T-\lambda I)^{m+1}. \]
\end{parts}
\end{exercise}

\begin{exercise}
Suppose $\F=\C$ and $S\in\Lin(V)$ is a unitary operator. Prove that the
constant term in the characteristic polynomial of $S$ has absolute value $1$.
\end{exercise}

\begin{exercise}
Suppose that $\F=\C$ and $V_1,\dots,V_m$ are nonzero subspaces of $V$ such
that
\[ V = V_1\oplus\cdots\oplus V_m. \]
Suppose $T\in\Lin(V)$ and each $V_k$ is invariant under $T$. For each $k$, let
$p_k$ denote the characteristic polynomial of $T|_{V_k}$. Prove that the
characteristic polynomial of $T$ equals $p_1\cdots p_m$.
\end{exercise}

\begin{exercise}
Suppose $p,q\in\Poly(\C)$ are monic polynomials with the same zeros and $q$ is
a polynomial multiple of $p$. Prove that there exists
$T\in\Lin(\C^{\deg q})$ such that the characteristic polynomial of $T$ is $q$
and the minimal polynomial of $T$ is $p$.
\exnote{This exercise implies that every monic polynomial is the characteristic
polynomial of some operator.}
\end{exercise}

\begin{exercise}
Suppose $A$ and $B$ are block diagonal matrices of the form
\[ A = \begin{pmatrix} A_1 & & 0\\ & \ddots & \\ 0 & & A_m \end{pmatrix},
   \quad
   B = \begin{pmatrix} B_1 & & 0\\ & \ddots & \\ 0 & & B_m \end{pmatrix}, \]
where $A_k$ and $B_k$ are square matrices of the same size for each
$k=1,\dots,m$. Show that $AB$ is a block diagonal matrix of the form
\[ AB = \begin{pmatrix} A_1B_1 & & 0\\ & \ddots & \\ 0 & & A_mB_m
   \end{pmatrix}. \]
\end{exercise}

\begin{exercise}
Suppose $\F=\R$, $T\in\Lin(V)$, and $\lambda\in\C$.
\begin{parts}
\item Show that $u+iv\in G(\lambda,T_{\C})$ if and only if
  $u-iv\in G(\overline{\lambda},T_{\C})$.
\item Show that the multiplicity of $\lambda$ as an eigenvalue of $T_{\C}$
  equals the multiplicity of $\overline{\lambda}$ as an eigenvalue of
  $T_{\C}$.
\item Use (b) and the result about the sum of the multiplicities
  (\ref{ax:8.25}) to show that if $\dim V$ is an odd number, then $T_{\C}$ has
  a real eigenvalue.
\item Use (c) and the result about real eigenvalues of $T_{\C}$ (Exercise 17 in
  Section 5A) to show that if $\dim V$ is an odd number, then $T$ has an
  eigenvalue (thus giving an alternative proof of \ref{ax:5.34}).
\end{parts}
\exnote{See Exercise 33 in Section 3B for the definition of the
complexification $T_{\C}$.}
\end{exercise}
\end{exc}
